\chapterhead{49}{ORR-10}{Sound Management of Risks Related to Money Laundering and Financing of Terrorism}{2}{2}

\lohead{49}{a}{Explain best practices recommended by the Basel Committee for the
assessment, management, mitigation, and monitoring of money laundering and
financing of terrorism (ML/FT) risks.}

\orsub{Management of money laundering and financial terrorism risks}

\orsubb{1) Risk assessment}

\begin{checks}
\item Banks need strong Customer Due Diligence (CDD) to prevent criminal
activity.
\item ML/FT risks must be assessed at four levels: country, sector, bank, and
individual customer relationships.
\item Banks need policies for
  \begin{lettered}
  \item Customer identification and due diligence
  \item Customer acceptance
  \item Monitoring business relationships and operations
  \end{lettered}
\item Understanding of ML/FT risks should consider
  \begin{lettered}
  \item Customer base
  \item Bank's products and services
  \item Delivery channels for products and services
  \item Jurisdictions of operation (bank and customers)
  \item Internal data (transactions) and external data (national risk
  assessments)
  \end{lettered}
\end{checks}

\orsubb{2) Risk management}

\begin{itemize}
\item \textbf{Governance:} the board of directors must oversee risk policies,
risk management activities, and compliance. ML/FT risk assessments need clear
communication to the board.
\item \textbf{Chief AML/CFT officer:} a qualified chief AML/CFT officer with
authority is required.
\end{itemize}

\orsubb{3) Risk mitigation}

\noindent\begin{minipage}{\textwidth}
\renewcommand{\arraystretch}{1.32}
\begin{tabularx}{\textwidth}{@{}YYY@{}}
\rowcolor{navy} \thd{1st Line: Business Units} & \thd{2nd Line: AML/CFT Officer} & \thd{3rd Line: Internal Audits}\\
{\tabfont
a)~Identify, assess, and control money laundering / financing of terrorism
(ML/FT) risks.\newline
b)~Written policies and procedures for employees.\newline
c)~Employee training on identifying and reporting suspicious transactions.\newline
d)~High ethical standards expected.} &
{\tabfont
a)~Monitors and fulfils AML/CFT duties.\newline
b)~Internal and external contact person for AML/CFT issues.\newline
c)~Reports directly to senior management / board.\newline
d)~Avoids conflicts of interest (no business line responsibilities).} &
{\tabfont
a)~Regularly audits the bank's AML/CFT policies.\newline
b)~Ensures effectiveness of AML/CFT controls.}\\
\end{tabularx}
\end{minipage}
\orfigcap{ML/FT risk mitigation across the three lines of defense.}

\orsubb{4) Risk monitoring}

Should match the bank's scale, activities, and complexity. Includes automated
systems for larger or international banks. If a bank chooses not to use IT
monitoring, it must document this decision and show that its alternative methods
are effective.

IT systems should enable accurate monitoring of customer transaction profiles and
help in filing Suspicious Transaction Reports (STRs) and mitigating ML/FT risks.

% ---------------------------------------------------------------------
\lohead{49}{b}{Describe recommended practices for the acceptance, verification,
and identification of customers at a bank.}

\orsub{Recommended practices for customer acceptance, verification, and
identification at banks}

\orsubb{1) Customer acceptance}

\begin{checks}
\item \textbf{Customer risk assessment:} assess risks based on the customer's
background, occupation, banking behaviour, and source of income and wealth.
\item \textbf{Follow a risk-based approach:} simplified procedures for low-risk
customers (e.g., low balances). Enhanced due diligence for high-risk customers:
  \begin{lettered}
  \item Large balances and frequent cross-border transfers
  \item Politically Exposed Persons (PEPs), and foreign PEPs
  \end{lettered}
\item \textbf{Customer acceptance standards:} not so restrictive as to deny
access to the general public. Banks must determine acceptable risk levels for
higher-risk customers.
\end{checks}

\orsubb{2) Customer verification}

\begin{checks}
\item Banks must identify and verify customer identities according to FATF
recommendations.
\item Establish a system for customer identification and verification.
\item In some cases, verify the identity of someone acting on behalf of the
owner(s).
\item Use reliable documents that are difficult to forge (e.g., passports, IDs).
\item Don't solely rely on written declarations of identity.
\item Implement procedures for cases where customer presence isn't possible,
especially for customers from high-risk jurisdictions.
\end{checks}

\orsubb{3) Customer identification}

\begin{checks}
\item \textbf{Risk profiles:} banks assess customer risk (transactions, purpose,
funds).
\item \textbf{Suspicious activity:} banks monitor for unusual transactions.
\item \textbf{CDD or STR:} banks do CDD before business, and file an STR
(Suspicious Transaction Report) if verification fails.
\item \textbf{High-risk / refused:} banks may refuse high-risk or anonymous
customers.
\item \textbf{Numbered accounts:} ID verification is still required for numbered
accounts.
\end{checks}

% ---------------------------------------------------------------------
\lohead{49}{c}{Explain practices for managing ML/FT risks in a group-wide and
crossborder context.}

\orsub{ML/FT risk management for cross-border banks}

\begin{lettered}
\item \textbf{Unified policies:} establish and apply uniform AML/CFT policies
globally, adjusting for stricter local laws as needed.
\item \textbf{Customer data sharing:} facilitate cross-border customer data
sharing within legal limits for risk management.
\item \textbf{Assessments and audits:} conduct comprehensive customer and group
risk assessments, plus audits.
\item \textbf{Business referrals:} apply own AML/CFT standards to referred
business, adjusting for stricter external standards when necessary.
\item \textbf{Customer monitoring:} consolidate monitoring of customers, owners,
and transactions.
\item \textbf{AML/CFT leadership:} designate a chief AML/CFT officer for global
compliance oversight.
\item \textbf{System centralisation:} centralise data systems in larger banks for
better risk management.
\end{lettered}

\orsub{Role of supervisors}

\begin{lettered}
\item Supervisors play a critical role in AML/CFT oversight by adhering to
international standards, setting supervisory expectations, and employing a
risk-based approach to monitor banks' ML/FT management.
\item They evaluate banks' risk assessments and the effectiveness of controls,
allocate resources for risk reviews, and ensure compliance with regulations.
\item Supervisors also work to harmonise jurisdictional requirements, verify
group-wide policy compliance, cooperate with home-country regulators, manage
group audits, maintain information confidentiality, and avoid being classified as
``third parties'' in restrictive jurisdictions.
\end{lettered}
